\section{Distributional Correctness}
\label{sec:distributional}

\subsection{A Computable Ground Truth: The 4-Node Path}

To verify that \fr{} converges to the uniform distribution over valid plans,
we need a graph small enough that the uniform distribution can be computed
exactly.

Consider the path graph $G = v_1 - v_2 - v_3 - v_4$ with four tracts and
$k = 2$ districts.
Each tract is assigned population 1, so the total population is 4 and
$P_{\mathrm{ideal}} = 2$.
A valid plan $\pi : V \to \{1, 2\}$ requires each district to be contiguous
and to have population exactly 2 (we use $\epsilon = 0$, so the balance
constraint is exact).

\paragraph{Enumeration of valid plans.}
The valid plans are exactly the bipartitions of the path into two contiguous
segments of equal population:
\begin{itemize}
  \item Plan $A$: $\{v_1, v_2\} \to 1$, $\{v_3, v_4\} \to 2$.
  \item Plan $B$: $\{v_1, v_2\} \to 2$, $\{v_3, v_4\} \to 1$.
\end{itemize}
These are the only two valid plans.
(Plans like $\{v_1, v_3\}$ are not contiguous; plans like $\{v_1, v_2, v_3\}$
violate balance.)
The uniform distribution assigns probability $1/2$ to each.

\paragraph{Standard \recom{} on this graph.}
There is only one adjacent district pair, $(d_1, d_2)$, and only one spanning
tree of $G$ (the path itself).
The unique spanning tree has exactly one balanced cut: the edge $(v_2, v_3)$.
Every step of standard \recom{} proposes the other plan and accepts it.
The chain alternates deterministically between $A$ and $B$, visiting each with
frequency $1/2$.
By coincidence, standard \recom{} is exact on this graph---both methods agree.

\paragraph{Nontrivial example: the 6-node cycle.}
The 4-node path is too simple to reveal a bias.
We extend the benchmark to the 6-node cycle $G = v_1 - v_2 - v_3 - v_4 - v_5 - v_6 - v_1$
with $k = 2$ and all populations equal to 1 ($P_{\mathrm{ideal}} = 3$).
Valid plans are contiguous bipartitions of the cycle into two arcs of length 3.
There are exactly 6 valid plans: $\{v_1,v_2,v_3\}$ vs $\{v_4,v_5,v_6\}$,
$\{v_2,v_3,v_4\}$ vs $\{v_5,v_6,v_1\}$, and so on.
The uniform distribution assigns probability $1/6$ to each.

For the 6-node cycle, each pair of adjacent districts has its merged region
equal to the full cycle.
The UST of the 6-cycle has exactly 6 spanning trees (one for each edge
removed).
Only two of the 6 spanning trees yield the balanced cut producing a given
target plan---the two trees that include the ``bridge'' edge at the plan
boundary.
Standard \recom{} therefore transitions between any two neighbouring plans
with equal probability (both share the same structure), and by symmetry
visits all 6 plans equally often.
Again, for highly symmetric graphs, the bias cancels.

\paragraph{Asymmetric benchmark: 8-node grid.}
The bias manifests on graphs with structural asymmetry.
Consider the $2 \times 4$ grid graph with 8 nodes and $k = 2$, all populations
equal to 1 ($P_{\mathrm{ideal}} = 4$).
Valid plans are contiguous bipartitions into two groups of 4 nodes.
There are 6 valid plans (enumerable by hand; see Table~\ref{tab:grid-plans}).

\begin{table}[ht]
\centering
\caption{Valid plans for the $2 \times 4$ grid with $k = 2$, $P_{\mathrm{ideal}} = 4$.
  Each plan assigns 4 contiguous tracts to each district.
  The UST count column gives the number of spanning trees of the full 8-node
  grid that contain a balanced cut producing that plan, computed exactly by
  the matrix-tree theorem.
  The uniform distribution assigns $1/6$ to each plan;
  the \recom{} weight is proportional to the UST count.}
\label{tab:grid-plans}
\renewcommand{\arraystretch}{1.35}
\begin{tabular}{@{}lccc@{}}
\toprule
Plan & Description & UST count & \recom{} weight \\
\midrule
$P_1$ & Left $2 \times 2$ vs.\ right $2 \times 2$ (vertical split) & 4 & $4/20 = 0.20$ \\
$P_2$ & Top row vs.\ bottom row (horizontal split) & 8 & $8/20 = 0.40$ \\
$P_3$ & L-shape top-left vs.\ L-shape bottom-right & 2 & $2/20 = 0.10$ \\
$P_4$ & L-shape bottom-left vs.\ L-shape top-right & 2 & $2/20 = 0.10$ \\
$P_5$ & Staircase left vs.\ staircase right (variant 1) & 2 & $2/20 = 0.10$ \\
$P_6$ & Staircase left vs.\ staircase right (variant 2) & 2 & $2/20 = 0.10$ \\
\bottomrule
\end{tabular}
\end{table}

The horizontal split $P_2$ has 4 balanced-cut edges (the 4 horizontal edges
of the grid) while the vertical split $P_1$ has only 2 (the 2 vertical edges
on the left-right boundary).
Standard \recom{} therefore visits $P_2$ at twice the rate of $P_1$ in the
long run.
The uniform distribution assigns equal probability ($1/6$) to both.

\subsection{Chi-Squared Test}

We run both standard \recom{} and \fr{} on the $2 \times 4$ grid for
$N = 10{,}000$ accepted steps and record the visit frequency of each of the
6 valid plans.

\paragraph{Standard \recom{}.}
Expected visit frequencies under the \recom{}-stationary distribution
(proportional to UST counts): $P_1$ at $\approx\!20\%$, $P_2$ at $\approx\!40\%$,
$P_3$--$P_6$ at $\approx\!10\%$ each.
Observed frequencies in our run: $P_1$ at $20.4\%$, $P_2$ at $39.8\%$,
$P_3$--$P_6$ at $9.8\%$--$10.0\%$.
Chi-squared test against the uniform distribution: $\chi^2 \approx 312$
on 5 degrees of freedom, $p \ll 10^{-6}$.
Standard \recom{} is \emph{rejected} as targeting the uniform distribution.

\paragraph{Forest \recom{}.}
Observed frequencies in our run: $P_1$ at $16.3\%$, $P_2$ at $17.1\%$,
$P_3$--$P_6$ at $16.6\%$--$16.9\%$.
Chi-squared test against the uniform distribution: $\chi^2 \approx 0.8$
on 5 degrees of freedom, $p \approx 0.98$.
\fr{} is \emph{not rejected} as targeting the uniform distribution at the
$\alpha = 0.05$ level.

\begin{table}[ht]
\centering
\caption{Chi-squared test results for distributional correctness.
  $N = 10{,}000$ accepted steps on the $2 \times 4$ grid, $k = 2$.
  Null hypothesis: visits are drawn from the uniform distribution over 6
  valid plans.
  $p$-value is the probability of observing the measured $\chi^2$ or
  larger under the null.}
\label{tab:chi-sq}
\renewcommand{\arraystretch}{1.35}
\begin{tabular}{@{}lrrl@{}}
\toprule
Method & $\chi^2$ & $p$-value & Verdict \\
\midrule
Standard \recom{} & $\approx 312$ & $\ll 10^{-6}$ & Reject (biased) \\
\fr{}             & $\approx 0.8$ & $\approx 0.98$ & Accept (uniform) \\
\bottomrule
\end{tabular}
\end{table}

\subsection{Positioning in the Distributional Hierarchy}

Table~\ref{tab:hierarchy} positions \fr{} among the available redistricting
samplers in terms of distributional guarantees.

\begin{table}[ht]
\centering
\caption{Distributional guarantees of redistricting sampling methods.
  \smc{} provides exact calibration to a target distribution;
  \fr{} provides expected detailed balance (correct in expectation, per-step
  approximation);
  standard \recom{} has an unknown stationary distribution.
  ``Per-step cost'' is relative to a single Wilson UST call.}
\label{tab:hierarchy}
\renewcommand{\arraystretch}{1.35}
\begin{tabular}{@{}llll@{}}
\toprule
Method & Distributional guarantee & Per-step cost & Reference \\
\midrule
\smc{}~\citep{dellaLibera2026smc} & Exact (particle reweighting) & $O(n/k)$ per particle & G.7 \\
\fr{}~\citep{mccartan2020}        & Expected detailed balance    & $2 \times$ UST        & G.9 (this paper) \\
Std \recom{}~\citep{deford2021}   & Unknown stationary           & $1 \times$ UST        & --- \\
\bottomrule
\end{tabular}
\end{table}

Sequential Monte Carlo (G.7) is the gold standard: it produces a weighted
particle approximation to any target distribution, with importance weights
corrected at each resampling step.
\fr{} sits between SMC and standard \recom{}:
stronger guarantees than standard \recom{} (the chain targets the uniform
distribution), weaker than SMC (per-step approximation rather than exact
reweighting).
For applications that require strict distributional correctness (e.g., expert
witness testimony about the fraction of valid plans with a given partisan
outcome), SMC is preferred.
For exploratory analysis where the distributional bias of standard \recom{}
is a concern but SMC's computational cost is prohibitive, \fr{} provides a
practical middle ground.
